\section{A finite cyclotomic generator for products of rational tails}
\label{cyc:sec}

The incoming Shifts and Twists report \cite{HTST} asks for an analogue of its two-parameter
alternating identity at higher roots of unity.  There is a precise answer
for one conjugation sector, valid at every conductor.  The answer below
uses only elementary trigonometric functions and finitely many digamma
values.  Its coefficient form is consistent with the classical
depth-two parity theorem for colored multiple polylogarithms
\cite{PanzerParity}.  We give an independent bilateral-sum proof, including
the limiting conventions at the trivial color.  No reduction of the
opposite conjugation sector is asserted.

\subsection{The finite functions}

If $\omega$ is a root of unity, define, initially for $|t|<1$,
\begin{equation}\label{cyc:F}
 F_\omega(t)=\begin{cases}
 -\psi(1-t)-\gamma,&\omega=1,\\[2pt]
 -\dfrac1N\displaystyle\sum_{j=1}^N\omega^j
                  \psi\bigl((j-t)/N\bigr),
       &\omega^N=1,\quad\omega\ne1.
 \end{cases}
\end{equation}
The value is independent of the chosen multiple $N$ of the conductor.
For $\omega\ne1$ it is the ordinary convergent series
$\sum_{k\ge1}\omega^k/(k-t)$; residue-class grouping and the digamma
difference series prove \eqref{cyc:F}.  Its Taylor series is
\[
 F_\omega(t)=\sum_{j\ge1}\operatorname{Li}_j(\omega)t^{j-1}.
\]
For the trivial color the coefficient of $t^0$ in this last notation is
set to zero; all coefficients of positive degree are their ordinary
zeta values.  This convention does not assign a value to a divergent
harmonic series.

Write $\omega=e^{2\pi i\theta}$ with $0<\theta<1$ for a nontrivial
color and put
\begin{equation}\label{cyc:B}
 B_\omega(t)=\begin{cases}
 \pi t\cot(\pi t),&\omega=1,\\[2pt]
 \dfrac{\pi t e^{i\pi(2\theta-1)t}}{\sin(\pi t)},&\omega\ne1.
 \end{cases}
\end{equation}
These functions have value $1$ at zero.  They obey
\begin{equation}\label{cyc:BF}
 B_\omega(t)=1-tF_\omega(t)+tF_{\omega^{-1}}(-t).
\end{equation}
For example $B_{-1}(t)=\pi t/\sin(\pi t)$.
The coefficients $b_j(\omega)=[t^j]B_\omega(t)$ are explicit:
\[
 b_j(\omega)=\frac{(2\pi i)^j}{j!}B_j(\theta)\quad(\omega\ne1),
 \qquad b_0(1)=1,\quad b_1(1)=0,
 \quad b_j(1)=-2\zeta(j)\ (j\ge2\hbox{ even}),
\]
and $b_j(1)=0$ for odd $j\ge3$.  The $B_j(\theta)$ here are ordinary
Bernoulli polynomials, whereas $B_\omega(t)$ denotes their scaled
generating function.  Formula \eqref{cyc:BF} also proves that for $j\ge1$
\[
 b_j(\omega)=-\operatorname{Li}_j(\omega)
                    -(-1)^j\operatorname{Li}_j(\omega^{-1}),
\]
with the stated convention at $\omega=1$, $j=1$.

Define the following finite expression:
\begin{align}\label{cyc:C}
 C_{a,b}(X,Y)={}&
 \frac{B_{ab}(X)F_b(Y-X)-B_a(X)F_b(Y)}X\notag\\
 &-\frac{B_{ab}(Y)F_{a^{-1}}(Y-X)-F_{a^{-1}}(-X)}Y.
\end{align}
Each numerator has the indicated factor, so $X=0$ and $Y=0$ are
removable for this expression near the origin.  We use the finite
four-point operation
\begin{equation}\label{cyc:Q}
 Q_{a,b}(x,y)=\frac{
 C_{a,b}(x,y)+C_{a,b}(x,-y)
 -C_{a,b}(-x,y)-C_{a,b}(-x,-y)}{4x}.
\end{equation}
Finally set
\begin{equation}\label{cyc:J}
 J_\omega(t)=\frac{B_\omega(t)-B_\omega(-t)}{2t}
 =\begin{cases}
 0,&\omega=1,\\[2pt]
 \displaystyle i\pi\frac{\sin(\pi(2\theta-1)t)}{\sin(\pi t)},
       &\omega\ne1.
 \end{cases}
\end{equation}
The function $J_\omega$ is even.  Thus
$[J_\omega(x)-J_\omega(y)]/(x^2-y^2)$ has its usual derivative
interpretation when $y=\pm x$.

\subsection{All colors and an arbitrary outer character}

For roots of unity $a,b,c$ write
\begin{align}
 E_n^a(x)&=\sum_{k>n}\frac{a^k}{k^2-x^2},\notag\\
 R_{a,b;c}(x,y)&=\sum_{n\ge0}c^n\{c(n+1)-n\}
                                      E_n^a(x)E_n^b(y).
 \label{cyc:R}
\end{align}
The normalization of the weight makes its partial sum exactly
$c^m m$.  In particular $R_{a,b;-1}=-\mathcal K_{a,b}$ for the
alternating weight $\mathcal K_{a,b}=\sum(-1)^n(2n+1)E_n^aE_n^b$.

\begin{theorem}[Cyclotomic rational-tail generator]\label{cyc:master}
Let $a,b,c$ be arbitrary roots of unity and put $d=abc$.  Initially
assume $|x|+|y|<1$ and $xy\ne0$.  Then
\begin{align}\label{cyc:mastereq}
 &R_{a,b;c}(x,y)+R_{a^{-1},b^{-1};c^{-1}}(x,y)\notag\\
 &\qquad=Q_{a,bc}(x,y)+Q_{b,ac}(y,x)
                 +\frac{J_d(x)-J_d(y)}{x^2-y^2}.
\end{align}
Both sides extend meromorphically to $\mathbb C^2$.  All poles away
from $x\in\mathbb Z\setminus\{0\}$ or
$y\in\mathbb Z\setminus\{0\}$ are removable in the combined
right side.  The series is absolutely locally normally convergent
unless $a=b=1$ and $c\ne1$.  In that exceptional case it converges
ordinarily and locally uniformly; subtracting the summand at
$(x,y)=(0,0)$ makes the remaining series absolutely locally normally
convergent.  When $x,y$ are real, the left side of
\eqref{cyc:mastereq} is $2\Re R_{a,b;c}(x,y)$.
\end{theorem}

In particular this is a conductor-four generator, with quarter-shift
digamma expressions furnished by \eqref{cyc:F}.  It also works for
odd conductors, because the outer color $c$ is arbitrary.  The appearance
of a conjugate pair is essential: the theorem proves a depth reduction
for the real sector at odd total weight, and does not prove such a
reduction for each complex summand separately.

For nonzero rational $x,y$ with $x,y,x+y,x-y\notin\mathbb Z$,
Gauss's rational digamma formula \cite[Eq.~(5.4.19)]{DLMFGamma}
makes the answer a finite explicit
combination of $\pi$, $\pi^2$, and $\pi\log\alpha$ for positive
algebraic $\alpha$, with algebraic coefficients.  The same assertion
holds on $x=y\ne0$ when $2x\notin\mathbb Z$.  To see the stated
coefficient class, $F_\omega$ at a rational argument is an algebraic
linear combination of $1$, $\pi$, and such logarithms: its Euler
constant cancels in \eqref{cyc:F}.  Every $B_\omega$ value at a
nonzero rational argument is $\pi$ times an algebraic number.
The apparently unmultiplied term $F_{a^{-1}}(-X)/Y$ in
\eqref{cyc:C} is odd in $Y$ and disappears in \eqref{cyc:Q}.
The remaining $J$ quotient is elementary, also at its diagonal
limit.  This statement concerns the displayed nonsingular rational
parameters; the axes can introduce higher polylogarithmic constants,
as the $\zeta(3)$ limits below demonstrate.

\subsection{A bilateral proof with finite cutoff control}

For nontrivial unit colors $a,b$, introduce the ordinary generating
series
\[
 Z_{a,b}(X,Y)=\sum_{k>m\ge1}\frac{a^k b^m}{(k-X)(m-Y)}.
\]
Summation by parts in $k$ proves locally uniform convergence: the inner
partial sums grow at most logarithmically and their successive
increments are $O(1/k)$.  We claim
\begin{equation}\label{cyc:bilateral-generator}
 Z_{a,b}(X,Y)-Z_{a^{-1},b^{-1}}(-X,-Y)
 =C_{a,b}(X,Y)-\frac{F_{ab}(X)-F_{ab}(Y)}{X-Y}.
\end{equation}
The last quotient is removable at $X=Y$.

To prove it, sum first over the full integer lattice in
\[
 S=\sum_{m\ge1}b^m\sum_{k\in\mathbb Z}
       \frac{(ab)^k}{(k-X)(k+m-Y)}.
\]
For each fixed $m$ the inner series is absolutely convergent.  The
bilateral partial-fraction formula is
\[
 \operatorname{PV}\sum_{k\in\mathbb Z}\frac{\omega^k}{k-t}
                  =-\frac{B_\omega(t)}t.
\]
It follows either from \eqref{cyc:BF} or from the elementary
Fourier series of the shifted resolvent.  The symmetric principal
value at $\omega=1$ is exactly the cotangent in \eqref{cyc:B}.
Partial fractions therefore give
\begin{equation}\label{cyc:S-eval}
 S=-\frac{B_{ab}(X)}X F_b(Y-X)
                  +\frac{B_{ab}(Y)}YF_{a^{-1}}(Y-X).
\end{equation}
Now put $l=k+m$.  The positive, negative, mixed, and two zero-index
parts of the region $l>k$ give respectively
\begin{align*}
 S={}&Z_{b,a}(Y,X)+Z_{a^{-1},b^{-1}}(-X,-Y)
       -F_{a^{-1}}(-X)F_b(Y)\\
    &-F_b(Y)/X+F_{a^{-1}}(-X)/Y.
\end{align*}
This rearrangement has a direct cutoff justification.  Keep
$1\le m\le M$ before separating regions.  In each same-sign region,
the omitted strip is bounded after summation by parts by a constant
times
\[
 \sum_{k\ge1}\frac1{k(k+M)}=O(\log M/M).
\]
In the mixed region the difference between the triangular cutoff
$u+v\le M$ and the product of the two ordinary sums is
$O(\log M/M)$ as well: for $u<M$ use the nontrivial-color tail
$O((M-u)^{-1})$, and for $u\ge M$ use the nontrivial-color tail in
$u$.  The two zero-index tails are $O(1/M)$.  These estimates are
uniform on compact parameter sets avoiding poles.  Thus all the
displayed limiting rearrangements are valid.

The elementary product decomposition is
\[
 Z_{b,a}(Y,X)=F_a(X)F_b(Y)-Z_{a,b}(X,Y)
                    -\frac{F_{ab}(X)-F_{ab}(Y)}{X-Y}.
\]
Use this identity, \eqref{cyc:BF}, and \eqref{cyc:S-eval} to obtain
\eqref{cyc:bilateral-generator}.

Here are explicit conventions for the boundary colors.  For $r\ge2$,
$s\ge1$, let $w=r+s$ and write $L_j(\omega)=\operatorname{Li}_j(\omega)$.
Coefficient extraction in \eqref{cyc:bilateral-generator} gives
\begin{align}\label{cyc:finite-parity}
 &Z(r,s;a,b)-(-1)^w Z(r,s;a^{-1},b^{-1})\notag\\
 ={}&(-1)^s\sum_{j=s}^{w}(-1)^j
          \binom{j-1}{s-1}L_j(b)b_{w-j}(ab)-L_w(ab)\notag\\
 &+(-1)^r\sum_{j=r}^{w}\binom{j-1}{r-1}
               L_j(a^{-1})b_{w-j}(ab)-L_s(b)b_r(a),
\end{align}
where
$Z(r,s;a,b)=\sum_{k>m\ge1}a^k b^m/(k^r m^s)$.
Every double series here is absolutely convergent.  If $b=1$, $s=1$,
the first summand with $j=1$ cancels the last term before any value
is assigned to $L_1(1)$.  If $ab=1$, the coefficient of $b_1(ab)$
vanishes, since $b=a^{-1}$ and
$\binom{w-2}{s-1}=\binom{w-2}{r-1}$ with opposite signs.
There is no other $L_1(1)$ term because $r\ge2$.  Consequently the
finite formula, and its normally convergent even-$r$, odd-$s$
generator, extend to every root-of-unity color with exactly
\eqref{cyc:F}--\eqref{cyc:B}.  One may also obtain this extension
by taking colors to $1$ on the unit circle; the only discarded
terms have $r=1$, which the projection in \eqref{cyc:Q} annihilates.

We now prove the tail theorem.  A finite outer cutoff, followed by
the telescoping identity
\[
 \sum_{n=0}^{m-1}c^n\{c(n+1)-n\}=c^m m,
\]
gives
\begin{equation}\label{cyc:ordered-tail}
 R_{a,b;c}(x,y)=D_{a,bc}(x,y)+D_{b,ac}(y,x)+H_{abc}(x,y),
\end{equation}
where
\begin{align*}
 D_{u,v}(x,y)&=\sum_{k>m\ge1}
   \frac{u^k v^m m}{(k^2-x^2)(m^2-y^2)},\\
 H_d(x,y)&=\frac{O_d(x)-O_d(y)}{x^2-y^2},\qquad
 O_d(t)=\frac{F_d(t)+F_d(-t)}2.
\end{align*}
The double sums are absolutely convergent, dominated by a constant
times $\sum_{k>m}k^{-2}m^{-1}$.  The omitted cutoff boundary is
$O(M^{-1})$, even when both tails are untwisted.  Applying the
four-point operation to \eqref{cyc:bilateral-generator} gives
\[
 D_{u,v}+D_{u^{-1},v^{-1}}=Q_{u,v}-H_{uv}.
\]
Add \eqref{cyc:ordered-tail} to its inverse-color version.  The two
copies of $-H_d$, together with the two diagonal terms, leave
$H_{d^{-1}}-H_d$.  Formula \eqref{cyc:BF} says
$O_{d^{-1}}-O_d=J_d$, proving \eqref{cyc:mastereq}.

For completeness, locally uniformly away from the coordinate poles,
\[
 E_n^1(x)=O(n^{-1}),\qquad E_n^a(x)=O(n^{-2})\quad(a\ne1),
 \qquad E_n^a(x)-E_n^a(0)=O(n^{-3}).
\]
These imply every asserted absolute convergence statement.  If
$a=b=1$ and $c\ne1$, the scalar summand at zero is
$c^n\{(c-1)/n+O(n^{-2})\}$ for $n\ge1$, which is ordinarily
convergent; the parameter difference is $O(n^{-3})$.  The case
$c=1$ has weight $1$ and is absolutely convergent even for $a=b=1$.
Near a coordinate pole only finitely many tail terms are singular.
Normal convergence and the meromorphic identity theorem therefore
give the claimed continuation and all oblique removability.

\subsection{Higher polynomial weights and their precise parity sector}

For $p,q\ge2$, $1\le h<\min(p,q)$, put
\[
 T_p^a(n)=\sum_{k>n}\frac{a^k}{k^p},\qquad
 M_{p,q,h}(a,b;c)=\sum_{n\ge0}c^n
       \{c(n+1)^h-n^h\}T_p^a(n)T_q^b(n).
\]
Exactly the same cutoff proof gives the ordinary identity
\begin{equation}\label{cyc:allweights}
 M_{p,q,h}(a,b;c)=Z(p,q-h;a,bc)+Z(q,p-h;b,ac)
                         +L_{p+q-h}(abc).
\end{equation}
The only possibly conditional case is $p=q=2$, $h=1$,
$a=b=1$, $c\ne1$; when $c=1$ the weight has one lower degree.
Let $w=p+q-h$.  Formula \eqref{cyc:finite-parity} now explicitly
reduces
\[
 M_{p,q,h}(a,b;c)-(-1)^w
 M_{p,q,h}(a^{-1},b^{-1};c^{-1})
\]
to finitely many single colored polylogarithm values and Bernoulli
coefficients.  For real parameters this is the real sector if $w$
is odd and the imaginary sector if $w$ is even.  Thus the parity
qualification is fully explicit at every conductor and every
admissible weight.

\subsection{Two conductor-four diagonal identities}

Write
\[
 O_-(x)=-\log2+O_+(x/2)-O_+(x),\qquad
 O_+(x)=-\frac{\psi(1+x)+\psi(1-x)+2\gamma}{2}.
\]
For real $x$ away from nonzero integers, define
\[
 D(x)=\Re\sum_{n\ge0}(-1)^n(2n+1)E_n^i(x)^2,\qquad
 V(x)=\sum_{n\ge0}(-1)^n(2n+1)|E_n^i(x)|^2.
\]
Both are ordinary absolutely convergent sums.  The all-color theorem
gives the compact evaluations
\begin{align}
 D(x)&=\frac{\pi}{2x\sin(\pi x)}
   \{\cos(\pi x/2)O_-(x/2)-\cos(\pi x)O_-(x)\},
 \label{cyc:Ddiag}\\
 V(x)&=\frac{\pi}{2x\sin(\pi x)}
   \{\cos(\pi x/2)O_-(x/2)-O_-(x)\}.
 \label{cyc:Vdiag}
\end{align}
One verification of the simplification is to use
\[
 O_i(t)+O_{-i}(t)=O_-(t/2),\quad
 \frac{B_i(t)+B_i(-t)}2
       =\frac{\pi t\cos(\pi t/2)}{\sin(\pi t)},
 \quad B_1(t)=\pi t\cot(\pi t)
\]
in \eqref{cyc:C}--\eqref{cyc:Q}.  At $a=b=i$ the product color in
the ordered sum is $1$, while at $a=i,b=-i$ it is $-1$;
in either case $J$ vanishes.  The terms at $F_{\pm i}(0)$ cancel
in the conjugate sum.  This gives \eqref{cyc:Ddiag} and
\eqref{cyc:Vdiag} directly, without a limiting singular quotient.

In particular, subtracting these two formulas proves the especially
short sine-tail square
\begin{equation}\label{cyc:sinesquare}
 \boxed{\displaystyle
 \sum_{n\ge0}(-1)^n(2n+1)
  \left(\sum_{k>n}\frac{\sin(\pi k/2)}{k^2-x^2}\right)^2
       =-\frac{\pi}{4x}\tan(\pi x/2)O_-(x).}
\end{equation}
The algebraic-square version extends meromorphically to complex $x$;
no complex conjugation is inserted there.  Its removable value at
zero is $\pi^2\log2/8$.  Two rational values are
\begin{align}
 \sum_{n\ge0}(-1)^n(2n+1)
  \left(\sum_{k>n}\frac{\sin(\pi k/2)}{k^2-1/4}\right)^2
       &=\frac\pi2,\label{cyc:sinehalf}\\
 \sum_{n\ge0}(-1)^n(2n+1)
  \left(\sum_{k>n}\frac{\sin(\pi k/2)}{k^2-1/16}\right)^2
       &=\pi(\sqrt2-1)
                 \{2-\sqrt2\log(1+\sqrt2)\}.
 \label{cyc:sinequarter}
\end{align}
Here $O_-(1/2)=-1$ and
$O_-(1/4)=\sqrt2\log(1+\sqrt2)-2$, by the rational digamma
formula.  The same substitution gives
\[
 D(1/2)=\pi\{\log(1+\sqrt2)-\sqrt2\},\qquad
 V(1/2)=\pi\{1+\log(1+\sqrt2)-\sqrt2\}.
\]
The coefficient limits are
\[
 D(0)=\frac9{32}\zeta(3)-\frac{3\pi^2\log2}{16},\qquad
 V(0)=\frac9{32}\zeta(3)+\frac{\pi^2\log2}{16}.
\]
For the cosine tails, the corresponding square has value
$9\zeta(3)/32-\pi^2\log2/16$ at zero.  These are evaluations
of an ordinary alternating outer sum, so a negative value of the
cosine square causes no positivity issue.

\subsection{An anchored quarter-color Gamma primitive}

Let $S(x)$ denote the holomorphic function in
\eqref{cyc:sinesquare} on $|x|<1$.  Then
\begin{equation}\label{cyc:gamma-primitive}
 \boxed{\displaystyle
 \int_0^u x\cot(\pi x/2)S(x)\,dx
  =\frac\pi8\log\!\left(
  \frac{\Gamma(1+u/2)\Gamma((1-u)/2)}
       {\Gamma(1-u/2)\Gamma((1+u)/2)}\right),\qquad |u|<1.}
\end{equation}
The integrand has a removable singularity at zero.  The logarithm
is its holomorphic branch equal to zero at $u=0$, and the path can
be any path in the unit disk.  Indeed \eqref{cyc:sinesquare} makes
the integrand $-\pi O_-(x)/4$.  Integrating the displayed digamma
formula for $O_-$ and then applying Gamma duplication yields
\eqref{cyc:gamma-primitive}; the two sides vanish at zero.
Absolute local normal convergence permits integration term by term
in the defining sum for $S$.  A concrete endpoint is
\[
 \int_0^{1/2}x\cot(\pi x/2)S(x)\,dx
       =\frac\pi4\log\frac{\Gamma(1/4)}{2\Gamma(3/4)}.
\]
This primitive is a direct conductor-four continuation of the
earlier alternating diagonal Gamma primitive, with its additive
constant fixed explicitly.
